\documentclass[11pt,a4paper]{article}
\usepackage[utf8]{inputenc}
\usepackage[T1]{fontenc}
\usepackage{geometry}
\geometry{margin=1in}
\usepackage{hyperref}
\usepackage{enumitem}
\usepackage{booktabs}
\usepackage{longtable}
\usepackage{xcolor}
\usepackage{titlesec}
\usepackage{amsmath,amssymb}

\titleformat{\section}{\Large\bfseries\color{blue!70!black}}{}{0em}{}
\titleformat{\subsection}{\large\bfseries\color{blue!50!black}}{}{0em}{}
\titleformat{\subsubsection}{\normalsize\bfseries\color{blue!30!black}}{}{0em}{}

\title{\textbf{Replication Plan}\\[0.5em]
\large Motion Tomography via Occupation Kernels\\[0.3em]
\normalsize OSTI ID: 1842593}
\author{Auto-generated Replication Blueprint}
\date{\today}

\begin{document}
\maketitle
\tableofcontents
\newpage

%---------------------------------------------------------------------
\section{Paper Overview}
%---------------------------------------------------------------------
This paper introduces a method for reconstructing spatially varying flow fields (e.g., ocean currents) from trajectory data of mobile sensors using \emph{occupation kernels} in a reproducing kernel Hilbert space (RKHS) framework. The key idea is that a trajectory passing through a region ``occupies'' that region, and the occupation kernel captures the integrated effect of the flow field along the trajectory. The method solves a linear system relating trajectory-integrated measurements to the underlying vector field, without requiring spatial discretization (meshless). The authors validate on synthetic 2D flow fields with known ground truth.

%---------------------------------------------------------------------
\section{Detailed Workflow}
%---------------------------------------------------------------------

\subsection{Phase 1: RKHS and Kernel Setup}
\begin{enumerate}[label=\textbf{1.\arabic*}]
  \item \textbf{Choose the base kernel function.}
    \begin{itemize}
      \item Gaussian RBF kernel: $k(x, y) = \exp(-\|x-y\|^2 / (2\sigma^2))$.
      \item Exponential kernel: $k(x, y) = \exp(-\|x-y\| / \ell)$.
      \item Set kernel bandwidth $\sigma$ or length scale $\ell$.
    \end{itemize}
  \item \textbf{Define the occupation kernel.}
    \begin{itemize}
      \item For a trajectory $\gamma: [0, T] \to \mathbb{R}^d$, the occupation kernel is:
      \begin{equation*}
        K_\gamma(x) = \int_0^T k(\gamma(t), x) \, dt
      \end{equation*}
      \item Compute via numerical quadrature (Simpson's rule or Gauss--Legendre).
    \end{itemize}
  \item \textbf{Compute the Gram matrix} of occupation kernels between all trajectory pairs:
    \begin{equation*}
      G_{ij} = \int_0^{T_i} \int_0^{T_j} k(\gamma_i(t), \gamma_j(s)) \, ds \, dt
    \end{equation*}
\end{enumerate}

\subsection{Phase 2: Synthetic Data Generation}
\begin{enumerate}[label=\textbf{2.\arabic*}]
  \item \textbf{Define ground-truth flow fields.}
    \begin{itemize}
      \item Double gyre flow: $\psi(x,y) = A \sin(\pi f(x,t)) \sin(\pi y)$ with time-dependent parameters.
      \item Rankine vortex, uniform flow with perturbation, or other specified fields.
    \end{itemize}
  \item \textbf{Generate synthetic trajectories.}
    \begin{itemize}
      \item Integrate the ODE $\dot{x} = v(x,t)$ using RK4 from various initial conditions.
      \item Add measurement noise to trajectory points if specified.
      \item Vary number of trajectories $N$ and trajectory duration $T$.
    \end{itemize}
\end{enumerate}

\subsection{Phase 3: Flow Field Reconstruction}
\begin{enumerate}[label=\textbf{3.\arabic*}]
  \item \textbf{Assemble the linear system.}
    \begin{itemize}
      \item The occupation kernel framework yields $\mathbf{G} \boldsymbol{\alpha} = \mathbf{b}$, where $\mathbf{b}$ encodes trajectory displacement data.
    \end{itemize}
  \item \textbf{Solve the regularized linear system} (Tikhonov regularization):
    \begin{equation*}
      \boldsymbol{\alpha} = (\mathbf{G} + \lambda \mathbf{I})^{-1} \mathbf{b}
    \end{equation*}
  \item \textbf{Evaluate the reconstructed field} at a grid of query points:
    \begin{equation*}
      \hat{v}(x) = \sum_i \alpha_i K_{\gamma_i}(x)
    \end{equation*}
\end{enumerate}

\subsection{Phase 4: Error Analysis and Validation}
\begin{enumerate}[label=\textbf{4.\arabic*}]
  \item \textbf{Compute reconstruction error.}
    \begin{itemize}
      \item RMSE between reconstructed and true flow field on a grid.
      \item Relative $L^2$ error norm.
    \end{itemize}
  \item \textbf{Convergence studies.}
    \begin{itemize}
      \item Error vs.\ number of trajectories.
      \item Error vs.\ kernel bandwidth.
      \item Error vs.\ regularization parameter $\lambda$.
    \end{itemize}
  \item \textbf{Reproduce paper figures} (flow field reconstructions, error maps, convergence plots).
\end{enumerate}

%---------------------------------------------------------------------
\section{Datasets}
%---------------------------------------------------------------------

\begin{longtable}{p{4.5cm} p{3.5cm} p{6cm}}
\toprule
\textbf{Dataset / Source} & \textbf{Type} & \textbf{Location / Notes} \\
\midrule
\endfirsthead
\toprule
\textbf{Dataset / Source} & \textbf{Type} & \textbf{Location / Notes} \\
\midrule
\endhead
Synthetic flow fields & Generated & Double gyre, vortex, etc.\ from analytic expressions in the paper \\
\midrule
Synthetic trajectories & Generated & ODE integration from specified initial conditions \\
\midrule
Paper figure data & Reference & For quantitative validation of reconstruction quality \\
\bottomrule
\end{longtable}

%---------------------------------------------------------------------
\section{Models, Tools, and Software}
%---------------------------------------------------------------------

\begin{longtable}{p{4.5cm} p{2.5cm} p{6.5cm}}
\toprule
\textbf{Tool / Library} & \textbf{Version} & \textbf{Purpose / URL} \\
\midrule
\endfirsthead
\toprule
\textbf{Tool / Library} & \textbf{Version} & \textbf{Purpose / URL} \\
\midrule
\endhead
Python & $\geq 3.8$ & Primary language \\
\midrule
NumPy & $\geq 1.21$ & Array operations, linear algebra \\
\midrule
SciPy & $\geq 1.7$ & ODE integration (\texttt{solve\_ivp}), quadrature, linear system solvers \\
\midrule
Matplotlib & $\geq 3.4$ & Flow field visualization, quiver plots, contour maps \\
\midrule
scikit-learn & $\geq 0.24$ & Optional: kernel computations, Gaussian process utilities \\
\bottomrule
\end{longtable}

\subsection{Hardware Requirements}
\begin{itemize}
  \item Any modern laptop. Gram matrix is $N \times N$ (number of trajectories), typically $N \sim 100$--$1000$.
\end{itemize}

\end{document}
